\chapter{Modele hierarchiczne: częściowe łączenie informacji między grupami}
\label{ch:hierarchical-bayes}
\paragraph{Powiązania i odsyłacze.} Wzory dla modeli sprzężonych wyprowadzono w rozdziale~\ref{ch:conjugate-families}; teorię ryzyka i decyzji w rozdziałach~\ref{ch:decision} i~\ref{ch:bayes-estimation-prediction}.


\section{Kiedy oddzielne analizy grup zawodzą?}
Załóżmy, że chcemy porównać niezawodność dziesięciu maszyn.
Dla jednej zarejestrowano dwa zdarzenia, a dla innej dwieście.
Estymowanie każdej intensywności zupełnie osobno pomija możliwość,
że maszyny są technologicznie podobne. Zastąpienie wszystkich
intensywności jedną wspólną wartością ukrywa natomiast rzeczywistą
zmienność pomiędzy maszynami. \emph{Model hierarchiczny} opisuje te
źródła niepewności osobno.

Niech $J\ge2$ oznacza liczbę grup, $j\in\{1,\ldots,J\}$ numer grupy,
$\theta_j$ jej nieznany parametr, a $Y_j$ jej obserwację zbiorczą.
Poniżej rozróżniamy trzy strategie: \emph{brak łączenia}
(ang. \emph{no pooling}), \emph{pełne łączenie}
(\emph{complete pooling}) oraz \emph{częściowe łączenie}
(\emph{partial pooling}). Trzecia strategia nie jest prostą średnią
wyników: jej wagi zależą od precyzji obserwacji i od zmienności
międzygrupowej.

\begin{definicja}[Hierarchiczny model normalny]
\label{def:normal-hierarchical}
Ustalmy $J\ge2$, znane liczby $v_j>0$, $\mu\in\R$ i $\tau^2>0$.
Niech $\theta_1,\ldots,\theta_J$ będą warunkowo niezależne przy
$(\mu,\tau^2)$ i $\theta_j\mid\mu,\tau^2\sim\Normal(\mu,\tau^2)$.
Warunkowo względem wektora $\theta=(\theta_1,\ldots,\theta_J)$
niech $Y_1,\ldots,Y_J$ będą niezależne oraz
$Y_j\mid\theta\sim\Normal(\theta_j,v_j)$.
Na przykład, gdy grupa $j$ zawiera $n_j\ge1$ niezależnych
pomiarów $X_{ji}\mid\theta_j\sim\Normal(\theta_j,\sigma^2)$
przy znanym $\sigma^2>0$, możemy przyjąć
$Y_j=n_j^{-1}\sum_{i=1}^{n_j}X_{ji}$ i $v_j=\sigma^2/n_j$.
Liczbę $\tau^2$ nazywamy \emph{wariancją międzygrupową}.
\end{definicja}

Zadanie praktyczne wymaga rozróżnienia poziomów losowości:
$\theta_j$ opisuje grupę, $X_{ji}$ pojedynczy pomiar, a $\mu$
wspólną lokalizację rozkładu parametrów grup. Jeśli $\mu,\tau^2$
są znane, poszczególne grupy pozostają niezależne marginalnie;
nie jest to jednak prawdą po dodaniu wspólnego losowego hiperparametru
$\mu$ i jego scałkowaniu.

\section{Dokładny posterior i pochodzenie częściowego łączenia}
\begin{twierdzenie}[Posterior normalny i współczynnik ściągania]
\label{thm:normal-partial-pooling}
W modelu z definicji~\ref{def:normal-hierarchical}, przy ustalonych
$\mu,\tau^2$, dla każdego $j$ i $y\in\R^J$ zachodzi
\[
\begin{aligned}
 \theta_j\mid (Y=y,\mu,\tau^2)&\sim\Normal(m_j,s_j^2),
 & s_j^2&=\left(\frac1{\tau^2}+\frac1{v_j}\right)^{-1},\\
 m_j&=w_jy_j+(1-w_j)\mu,
 & w_j&=\frac{\tau^2}{\tau^2+v_j}.
\end{aligned}
\]
Współrzędne $\theta_j$ są także niezależne warunkowo względem
$Y=y,\mu,\tau^2$. Przy $v_j=\sigma^2/n_j$ waga $w_j$ rośnie
wraz z liczebnością $n_j$ i wraz ze zmiennością $\tau^2$.
\end{twierdzenie}
\begin{proof}
Warunkowa gęstość $\theta_j=t$ po zaobserwowaniu $Y_j=y_j$
jest proporcjonalna do
$\exp[-(t-\mu)^2/(2\tau^2)-(y_j-t)^2/(2v_j)]$.
Współczynnik przy $t^2$ po rozwinięciu wynosi
$\tau^{-2}+v_j^{-1}$, a liniowy
$\mu\tau^{-2}+y_jv_j^{-1}$. Uzupełnienie kwadratu
daje wariancję $s_j^2$ i średnią
$s_j^2(\mu/\tau^2+y_j/v_j)=m_j$.
Iloczynowy charakter prioru i wiarygodności dowodzi
warunkowej niezależności po aktualizacji. Zapis
$w_j=n_j\tau^2/(n_j\tau^2+\sigma^2)$ wyjaśnia monotoniczność.
\end{proof}

\begin{twierdzenie}[Zysk w średnim błędzie i jego ograniczenie]
\label{thm:shrinkage-bayes-risk}
W założeniach twierdzenia~\ref{thm:normal-partial-pooling}
porównajmy estymatory $\delta_{\rm odd}(Y_j)=Y_j$,
$\delta_{\rm wsp}(Y_j)=\mu$ i
$\delta_{\rm hier}(Y_j)=m_j$ przy stracie kwadratowej.
Ich ryzyka uśrednione względem prioru $\theta_j\sim\Normal(\mu,\tau^2)$
wynoszą odpowiednio $v_j$, $\tau^2$ oraz
$\tau^2v_j/(\tau^2+v_j)$.
Ostatnia liczba jest ściśle mniejsza od obu poprzednich.
Przy \emph{ustalonym} $\theta_j=t$ ryzyko częstościowe estymatora
hierarchicznego wynosi jednak
$(1-w_j)^2(t-\mu)^2+w_j^2v_j$ i może być większe od $v_j$.
\end{twierdzenie}
\begin{proof}
Estymator $Y_j$ ma błąd $Y_j-\theta_j$ o wariancji $v_j$.
Estymator stały $\mu$ ma średni błąd kwadratowy
$\E[(\theta_j-\mu)^2]=\tau^2$.
Z twierdzenia~\ref{thm:normal-partial-pooling} posteriorowy
błąd średniokwadratowy $m_j$ równa się $s_j^2$ niezależnie od
obserwacji, zatem prawo całkowitej wartości oczekiwanej
daje ryzyko Bayesa $s_j^2$.
Alternatywnie, dla ustalonego $t$ podstawiamy
$m_j-t=w_j(Y_j-t)+(1-w_j)(\mu-t)$ i wykorzystujemy
zerową średnią szumu oraz $\Var(Y_j\mid\theta_j=t)=v_j$.
Uśrednienie ostatniego wzoru po $t$ daje
$w_j^2v_j+(1-w_j)^2\tau^2=s_j^2$.
Z uwagi na $v_j,\tau^2>0$ mamy
$s_j^2<v_j$ i $s_j^2<\tau^2$. Ryzyko warunkowe rośnie
natomiast kwadratowo z $|t-\mu|$, co wyklucza dominację
przy wszystkich $t$ w tym jednowymiarowym porównaniu.
\end{proof}

\begin{przyklad}[Dwie grupy o różnej liczebności]
Niech $\mu=10$, $\tau^2=4$, $\sigma^2=16$.
Grupa A ma $n_A=4$ i średnią $Y_A=14$, grupa B ma
$n_B=64$ i średnią $Y_B=14$. W pierwszej grupie
$v_A=4$, $w_A=1/2$, $m_A=12$, $s_A^2=2$.
W drugiej $v_B=1/4$, $w_B=16/17$, $m_B=234/17\approx13{,}765$,
$s_B^2=4/17\approx0{,}235$.
Obie grupy mają tę samą średnią z próby, ale wynikiem
analizy hierarchicznej jest silniejsze ściągnięcie grupy
słabo reprezentowanej i znacznie mniejsza niepewność
dla grupy z wieloma pomiarami. Nie oznacza to automatycznie,
że dane z grupy A są błędne.
\end{przyklad}


\Needspace{92mm}
\begin{center}
\begin{tikzpicture}[>=Latex, every node/.style={font=\small},node distance=12mm]
\node[draw,rounded corners,fill=blue!5,minimum width=30mm,minimum height=8mm] (hyper) {$\mu,\tau^2$};
\node[draw,rounded corners,below left=13mm and 18mm of hyper,minimum width=25mm] (t1) {$\theta_1$};
\node[draw,rounded corners,below right=13mm and 18mm of hyper,minimum width=25mm] (t2) {$\theta_J$};
\node[draw,rounded corners,below=11mm of t1] (y1) {$Y_1$};
\node[draw,rounded corners,below=11mm of t2] (y2) {$Y_J$};
\draw[->] (hyper)--(t1); \draw[->] (hyper)--(t2);
\draw[->] (t1)--(y1); \draw[->] (t2)--(y2);
\node at ($(t1)!0.5!(t2)$) {$\cdots$};
\node[below=2mm of y1,font=\footnotesize] {$v_1$ -- znane};
\node[below=2mm of y2,font=\footnotesize] {$v_J$ -- znane};
\end{tikzpicture}

{\small\emph{Rysunek 21.1.} Wspólne hiperparametry wyznaczają rozkłady parametrów grup, a obserwacje informują o poszczególnych grupach.}
\end{center}

\section{Prognoza dla znanej i nowej grupy}
\begin{twierdzenie}[Dwie różne predykcje]
\label{thm:hierarchical-predictive}
W modelu~\ref{def:normal-hierarchical} niech przyszły pomiar
$Z_j\mid\theta_j\sim\Normal(\theta_j,\sigma^2)$ będzie warunkowo
niezależny od dotychczasowych danych dla znanej grupy $j$,
a $Z_{\rm nowy}\mid\theta_{\rm nowy}\sim\Normal(\theta_{\rm nowy},\sigma^2)$
dla nowej grupy, gdzie $\theta_{\rm nowy}\mid\mu,\tau^2\sim\Normal(\mu,\tau^2)$
i jest niezależny od istniejących grup. Wtedy
\[
 Z_j\mid Y=y,\mu,\tau^2\sim\Normal(m_j,\sigma^2+s_j^2),
 \qquad Z_{\rm nowy}\mid Y=y,\mu,\tau^2\sim\Normal(\mu,\sigma^2+\tau^2).
\]
\end{twierdzenie}
\begin{proof}
Nowy pomiar jest sumą $\theta_j$ i niezależnego
szumu $\Normal(0,\sigma^2)$. Po aktualizacji $\theta_j$
ma rozkład normalny ze średnią $m_j$ i wariancją $s_j^2$,
co daje pierwszy rozkład. Dla nowej grupy obserwacje
$Y$ nie aktualizują jej indywidualnego parametru przy
znanych hiperparametrach: pozostaje rozkład
$\Normal(\mu,\tau^2)$. Splot z szumem daje drugą tezę.
\end{proof}

\begin{uwaga}[Jednostka walidacji]
Jeżeli prognozujemy kolejny pomiar \emph{tej samej} maszyny,
możemy trenować na jej wcześniejszych wynikach. Jeśli chcemy
przewidywać wyniki \emph{zupełnie nowej} maszyny, walidacja
powinna odkładać całe maszyny (\emph{leave-one-group-out}).
Losowy podział pojedynczych pomiarów między trening i test
zwykle przecenia jakość przewidywania dla nowych grup.
\end{uwaga}

\section{Nieznane hiperparametry: pełny Bayes i empirical Bayes}
W praktyce $\mu$ i $\tau^2$ nie są znane. Dotychczasowe
wzory należy wtedy \emph{uśrednić po ich posteriorze},
a nie podstawiać oszacowania i udawać, że jest pewne.

\begin{twierdzenie}[Posterior hiperparametrów w modelu normalnym]
\label{thm:normal-hyperposterior}
Przyjmijmy model~\ref{def:normal-hierarchical} z nieznanymi
hiperparametrami $\mu\in\R$, $u=\tau^2>0$, znanymi
$v_j>0$ oraz niezależnymi, właściwymi priorami
$\mu\sim\Normal(m_0,V_0)$ i
$u\sim\operatorname{InvGamma}(a,b)$, gdzie $V_0,a,b>0$.
Przez $\operatorname{InvGamma}(a,b)$ rozumiemy rozkład
z gęstością $b^a u^{-a-1}e^{-b/u}/\Gamma(a)$ dla $u>0$
względem $\lambda_1$. Po scałkowaniu parametrów grup
$Y_j\mid\mu,u\sim\Normal(\mu,u+v_j)$ są niezależne,
a posterior hiperparametrów ma gęstość proporcjonalną do
\[
 \exp\!\left[-\frac{(\mu-m_0)^2}{2V_0}\right]
 u^{-a-1}e^{-b/u}\prod_{j=1}^J(u+v_j)^{-1/2}
 \exp\!\left[-\frac{(y_j-\mu)^2}{2(u+v_j)}\right].
\]
Dla ustalonego $u$ posterior $\mu\mid u,Y=y$ jest normalny
o wariancji $V(u)=(V_0^{-1}+\sum_j(u+v_j)^{-1})^{-1}$
i średniej
$M(u)=V(u)(m_0/V_0+\sum_jy_j/(u+v_j))$.
Posterior hiperparametrów jest właściwy dla każdego $y\in\R^J$.
\end{twierdzenie}
\begin{proof}
Po scałkowaniu niezależnych $\theta_j$ każda obserwacja
$Y_j$ jest sumą dwóch niezależnych zmiennych normalnych,
stąd wynika jej gęstość i niezależność warunkowa.
Mnożąc iloczyn tych gęstości przez priory otrzymujemy
pierwszy wzór. Jako funkcja $\mu$, wykładnik jest
ujemną połową dodatnio określonej funkcji kwadratowej;
uzupełnienie kwadratu daje $M(u)$ i $V(u)$.
Ponieważ $u+v_j\ge v_j>0$, marginalna wiarygodność
$\prod_j\varphi(y_j;\mu,u+v_j)$ jest ograniczona
z góry przez $\prod_j(2\pi v_j)^{-1/2}$.
Jej całka względem właściwego prioru jest skończona
i dodatnia, co dowodzi normalizowalności posterioru.
\end{proof}

\begin{stwierdzenie}[Dodatkowa niepewność hiperparametrów]
\label{prop:hier-variance-hyper}
Niech $H=(\mu,u)$ ma właściwy posterior z poprzedniego
twierdzenia i niech drugi moment $\theta_j\mid Y$ będzie skończony.
Wtedy
\[
 \E[\theta_j\mid Y]=\E[m_j(H)\mid Y],\qquad
 \Var(\theta_j\mid Y)=\E[s_j^2(H)\mid Y]
                       +\Var(m_j(H)\mid Y).
\]
Tu $m_j(H),s_j^2(H)$ otrzymujemy ze wzorów
z twierdzenia~\ref{thm:normal-partial-pooling} po podstawieniu
$\tau^2=u$. Ostatni składnik jest nieujemny i znika tylko,
gdy posteriorowa średnia grupy nie zależy losowo od $H$.
\end{stwierdzenie}
\begin{proof}
Pierwszy wzór to własność iterowanej warunkowej wartości
oczekiwanej, drugi to prawo całkowitej wariancji
z warunkowaniem najpierw po $H$, a następnie po $Y$.
\end{proof}

\begin{uwaga}[Kiedy wersja empirical Bayes może mylić?]
Metoda \emph{empirical Bayes} szacuje $H$ z danych, a potem
podstawia $\widehat H$ do wzorów dla $m_j$ i $s_j^2$.
Bywa szybka i użyteczna przy wielu grupach, ale naiwny
przedział $m_j(\widehat H)\pm1{,}96s_j(\widehat H)$
nie obejmuje niepewności oszacowania $H$. Pełny model
bayesowski ją uwzględnia, choć rezultat nadal zależy od
modelu i hiperprioru. Przy kilku grupach $u$ jest często
słabo zidentyfikowany; trzeba obejrzeć jego posterior,
a nie tylko wybraną estymatę punktową.
\end{uwaga}

\section{Zliczenia i udziały: hierarchiczny model beta--dwumianowy}
\begin{twierdzenie}[Łączenie informacji o rzadkich wadach]
\label{thm:hierarchical-beta-binomial}
Niech $J\ge2$, $\kappa>0$, $m\in(0,1)$.
Warunkowo przy znanych $m,\kappa$ zmienne $p_j$ są niezależne
i $p_j\sim\operatorname{Beta}(\kappa m,\kappa(1-m))$.
Dla każdego $j$ niech $S_j\mid p_j\sim\Bin(n_j,p_j)$,
$n_j\ge1$, i niech zliczenia są warunkowo niezależne.
Wtedy dla $s_j\in\{0,\ldots,n_j\}$
\[
 p_j\mid S_j=s_j,m,\kappa
 \sim\operatorname{Beta}(\kappa m+s_j,\kappa(1-m)+n_j-s_j),
\]
\[
 \E[p_j\mid S_j=s_j,m,\kappa]
 =\frac{n_j}{n_j+\kappa}\frac{s_j}{n_j}
 +\frac{\kappa}{n_j+\kappa}m.
\]
Jeżeli w jednej grupie obserwujemy wskaźniki
$X_{j1},\ldots,X_{jn_j}$ niezależne warunkowo przy $p_j$
i o rozkładzie $\operatorname{Bernoulli}(p_j)$, to dla $i\ne k$
\[
 \Cov(X_{ji},X_{jk})=\frac{m(1-m)}{\kappa+1},\quad
 \Var(S_j)=n_jm(1-m)\left(1+\frac{n_j-1}{\kappa+1}\right).
\]
\end{twierdzenie}
\begin{proof}
Iloczyn wiarygodności dwumianowej i gęstości beta
jest proporcjonalny do
$p_j^{\kappa m+s_j-1}(1-p_j)^{\kappa(1-m)+n_j-s_j-1}$,
co daje pierwszy wzór i średnią. Dla różnych obserwacji
tej samej grupy, przy warunkowej niezależności,
$\E[X_{ji}X_{jk}]=\E[p_j^2]$ oraz
$\E X_{ji}=\E p_j=m$.
Wobec tego kowariancja wynosi
$\Var(p_j)=m(1-m)/(\kappa+1)$.
Ponieważ marginalnie $X_{ji}$ jest Bernoulliego
z parametrem $m$, suma $n_j$ wariancji i
$n_j(n_j-1)$ kowariancji daje ostatni wzór.
\end{proof}

\begin{przyklad}[Dwie linie produkcyjne]
Przyjmijmy $m=0{,}10$, $\kappa=20$. Jedna linia nie wykazała
wad w dwóch próbach, druga wykazała siedem wad w stu.
Oddzielne estymatory MLE to odpowiednio $0$ i $0{,}07$.
Posteriorowe średnie są równe
$(0+2)/(2+20)=0{,}0909\ldots$ oraz
$(7+2)/(100+20)=0{,}075$.
Pierwsza linia nie jest uznawana za całkowicie bezawaryjną
tylko na podstawie dwóch obserwacji. Przy niewłaściwie
ustalonym $m$ ten efekt może jednak prowadzić do błędnego
ściągania, więc parametry populacyjne należy poddawać
diagnostyce i analizie wrażliwości.
\end{przyklad}

\section{Zaskakujący rezultat: ściąganie może pomagać bez prioru}
Poniższe twierdzenie nie zakłada, że nieznany wektor średnich
został wylosowany z rozkładu. Stanowi częstościowy argument,
dlaczego estymowanie grup w pełnej izolacji może być nieefektywne.

\begin{twierdzenie}[Fenomen Jamesa--Steina]
\label{thm:james-stein}
Niech $d\ge3$, $\sigma^2>0$ i
$X\sim\Normal_d(\theta,\sigma^2 I_d)$, gdzie
$\theta\in\R^d$ jest ustalonym, nieznanym parametrem,
a $\Normal_d$ oznacza $d$-wymiarowy rozkład normalny.
Przy stracie $L(\theta,a)=\norm{a-\theta}^2$ estymator
$\delta_{\rm JS}(x)=(1-(d-2)\sigma^2/\norm{x}^2)x$
dla $x\ne0$ oraz $\delta_{\rm JS}(0)=0$ ma ryzyko
\[
 R(\theta,\delta_{\rm JS})
 =d\sigma^2-(d-2)^2\sigma^4
   \E_\theta\!\left[\frac1{\norm{X}^2}\right]
 <d\sigma^2=R(\theta,X)
\]
dla \emph{każdego} $\theta\in\R^d$.
Wartość oczekiwana we wzorze jest skończona i dodatnia.
\end{twierdzenie}
\begin{proof}
Niech $a=(d-2)\sigma^2$ i
$g(x)=-ax/\norm{x}^2$ dla $x\ne0$.
W przestrzeni $\R^d\setminus\{0\}$ zachodzi
$\norm{g(x)}^2=a^2/\norm{x}^2$ oraz
$\operatorname{div}g(x)=-a(d-2)/\norm{x}^2$.
Dla gęstości normalnej $f_\theta$ mamy
$\partial_i f_\theta(x)=-(x_i-\theta_i)f_\theta(x)/\sigma^2$.
Całkowanie przez części po obszarze
$\{\varepsilon<\norm{x}<R\}$, a następnie przejście
$\varepsilon\downarrow0$, $R\uparrow\infty$, daje
\[
 \E_\theta[(X-\theta)\cdot g(X)]
   =\sigma^2\E_\theta[\operatorname{div}g(X)].
\]
Istotnie, wyraz brzegowy na sferze o promieniu
$\varepsilon$ jest rzędu $O(\varepsilon^{d-2})$,
ponieważ $\norm{g(x)}=a/\varepsilon$, pole sfery ma
rząd $\varepsilon^{d-1}$, a $f_\theta$ jest lokalnie
ograniczona. Zewnętrzny wyraz brzegowy zanika dzięki
gaussowskiemu zanikowi gęstości.
Całkowalność $\norm{X}^{-2}$ w otoczeniu zera
wynika z $\int_0^1r^{d-3}\,\dd\lambda_1(r)<\infty$,
co wymaga $d\ge3$; poza kulą jednostkową funkcja
$\norm{x}^{-2}$ jest ograniczona.
Rozwijając błąd kwadratowy $X+g(X)-\theta$ dostajemy
\[
 R(\theta,X+g)=d\sigma^2
   +2\sigma^2\E_\theta[\operatorname{div}g(X)]
   +\E_\theta\norm{g(X)}^2.
\]
Podstawienie wyliczonych pochodnych daje
$d\sigma^2+[a^2-2\sigma^2a(d-2)]\E_\theta[\norm{X}^{-2}]$.
Dla $a=(d-2)\sigma^2$ współczynnik wynosi
$-(d-2)^2\sigma^4<0$. Ponieważ rozkład normalny nie
przypisuje dodatniej masy punktowi zero, definicja w zerze
nie wpływa na ryzyko.
\end{proof}

\begin{uwaga}[Nie należy nadinterpretować twierdzenia]
Wynik nie mówi, że wszystkie współrzędne osobno są
estymowane lepiej, ani że procedura Jamesa--Steina zawsze
zwraca wartość między zerem a średnią z próby (czynnik
ściągania może być ujemny). Stwierdza poprawę łącznego
częstościowego ryzyka kwadratowego dla $d\ge3$.
Łączenie informacji między grupami może więc przynieść
korzyść nawet bez bayesowskiej interpretacji $\theta$.
\end{uwaga}

\section{Kiedy jaką metodę stosować?}
\begin{center}\small
\begin{tabular}{>{\raggedright\arraybackslash}p{.27\textwidth}>{\raggedright\arraybackslash}p{.31\textwidth}>{\raggedright\arraybackslash}p{.27\textwidth}}
\toprule
\textbf{Sytuacja} & \textbf{Metoda do rozważenia} & \textbf{Główne ryzyko}\\
\midrule
Dużo obserwacji w każdej odrębnej grupie & Estymacja oddzielna lub hierarchiczna & Nadmierne ściąganie rzeczywiście różnych grup.\\
Małe próby w wielu podobnych grupach & Model hierarchiczny z partial pooling & Nietrafny rozkład grup i hiperprior.\\
Jedna wspólna intensywność uzasadniona procesem & Model ze wspólnym parametrem & Ukrycie heterogeniczności.\\
Dane binarne w licznych partiach & Beta--dwumianowy lub logit hierarchiczny & Zależności i zmiany w czasie.\\
Szybka analiza wielu grup & Empirical Bayes & Niedoszacowanie niepewności hiperparametrów.\\
Predykcja nowej grupy & Pełny posterior predykcyjny grupy & Niepoprawny podział trening/test.\\
\bottomrule
\end{tabular}
\end{center}

\Needspace{23\baselineskip}
\begin{lstlisting}[language=Python,caption={Kompletny program: partial pooling dla średnich grup i ryzyko Jamesa--Steina (Python, NumPy).}]
import numpy as np

mu, tau2, sigma2 = 10.0, 4.0, 16.0
n = np.array([4, 64], dtype=float)
y = np.array([14.0, 14.0])
v = sigma2 / n
w = tau2 / (tau2 + v)
post_mean = w * y + (1.0 - w) * mu
post_var = tau2 * v / (tau2 + v)
print("group posterior means:", post_mean)
print("group posterior SD:", np.sqrt(post_var))

rng = np.random.default_rng(2026)
d, sigma, repetitions = 5, 2.0, 100000
theta = np.zeros(d)
X = rng.normal(theta, sigma, size=(repetitions, d))
r2 = np.sum(X**2, axis=1)
js = (1.0 - (d - 2) * sigma**2 / r2)[:, None] * X
risk_raw = np.mean(np.sum((X - theta)**2, axis=1))
risk_js = np.mean(np.sum((js - theta)**2, axis=1))
print("empirical risks raw/JS:", risk_raw, risk_js)
print("theoretical risk raw:", d * sigma**2)
\end{lstlisting}

\begin{zadanie}[Rola precyzji i zmienności międzygrupowej]
Przy znanym $\mu=0$ i $v_j=1$ dwie analizy używają
odpowiednio $\tau^2=1/9$ oraz $\tau^2=9$.
Obliczyć $w_j$ i wyjaśnić różnicę.
\par\emph{Rozwiązanie.} Otrzymujemy $w_j=1/10$ i $9/10$.
W pierwszym przypadku podobieństwo grup jest silnym
założeniem priorowym i szacowanie ściąga prawie całkowicie
do zera. W drugim obserwowana grupa może wyraźnie
różnić się od populacji, więc wynik jest bliski średniej
z jej własnych pomiarów. Przed wyborem $\tau^2$ trzeba
sprawdzić zakres możliwej heterogeniczności i liczbę grup.
\end{zadanie}

% ================================================================
% ROZDZIAL 22 -- ASYMPTOTYKA BAYESOWSKA
% ================================================================
